\scanpage{55}
et est remplacé par son produit par l'antipodisme
(à droite, ou à gauche, c'est kif kif) si on remplace
$\Phi$ par $\Phi'$.

On appellera \uline{biicosaèdre orienté} une \deleted{système}
\deleted{famille à deux éléments $\{I,I'\}$ de \unsure{la une}}
\deleted{\unsure{str|de}} un polyèdre comb. orienté avec
deux comp. connexes $I,I'$, qui sont
des icosaèdres, et la donnée d'un ens. \inserted{$\{\theta,\theta'\}$}
de \uline{deux} bijections \inserted{(= correspondances)} entre $S(\deleted{\unsure{J}}I)$ et $S(I')$
\deleted{\unsure{chacun de la|classe de}}
commutant aux antipodismes, et
\deleted{\unsure{formant}} tel que $\theta'$ et $\theta$ sont \unsure{déduits|déduites|déduisent} l'un de
l'autre par comp. avec l'antipodisme,
et tel que, en passant au quotient \deleted{\unsure{sous}}\inserted{de $\theta$}
\unsure{et|ou} $\theta'$ modulo antipodisme, la
bijection qu'on obtient entre $S(I)/\underline a$ et $S(I')/\underline a$
\deleted{transforme la bijection de \unsure{icos|l'icos}}
donne sur l'ens commun deux structures
d'icosaèdres gauches \deleted{\unsure{opp}} complémentaires!
% slanted
\margin{et les bijections, \unsure{antér.ieures|antérieurement|extérieures} soient définies \unsure{resp}$^t$ par $\varepsilon^S_{\Phi,\Phi'}$ et $\varepsilon^S_{\Phi',\Phi}$}

Il revient au $=$ de dire que l'on se donne
\deleted{une donnée d'ident. antérieure des}
\deleted{$I'$ avec l'icos. complémentaire à}
\inserted{une paire d'iso orientés}
\[
\left\{\begin{aligned}
\overline I'&\xrightarrow[\sim]{\alpha_{I'}}\operatorname{compl}.(\overline I)\\
\overline I&\xrightarrow[\sim]{\alpha_I}\operatorname{compl}(\overline I')
\end{aligned}\right.
\]
\margin{icos. ``complémentaire'' à $\overline I$, avec l'orientation canonique déduite}
dont les composés soient les antipodismes
\[
\unsure{\alpha_I}\ \operatorname{compl}(\alpha_{I'})\alpha_I:
I\mathrel{\unsure{-|\to}}\underbrace{\operatorname{compl}\operatorname{compl}(I)}_{
\substack{\Vert\\-I,\ \text{i.e. }I\\\text{avec l'orientation}\\\text{opposée}}}
\]
soit l'antipodisme, qui est un
iso de $I$ avec $-I$
% archive sheet number: 54
